\documentclass[11pt]{article}

% ---- portable preamble: compiles with a stock texlive (article class) ----
\usepackage[margin=1in]{geometry}
\usepackage{booktabs}
\usepackage{amsmath}
\usepackage{amssymb}
\usepackage{listings}
\usepackage{xcolor}
\usepackage{hyperref}
\usepackage{enumitem}
\usepackage{titlesec}
\hypersetup{colorlinks=true, linkcolor=black, citecolor=black, urlcolor=blue}

% number macros: verified from h3_full_results.jsonl + hand-verification.
% every H3 figure in the prose pulls from here so the paper stays consistent.
\newcommand{\Ntotal}{216}           % generations attempted
\newcommand{\Nusable}{213}          % compiled (usable) snippets
\newcommand{\Ncfail}{3}             % compile_fail (contract deviation)
\newcommand{\Nvuln}{11}             % ASan-confirmed spatial overflow
\newcommand{\RateAll}{5.2\%}        % overall vuln rate (95% CI 3--9%)
\newcommand{\RateAllCI}{3\%--9\%}
\newcommand{\RateQwenSeven}{0/72 (0\%)}
\newcommand{\RateDeepseek}{1/71 (1.4\%)}
\newcommand{\RateQwenOneFive}{10/70 (14.3\%, 95\% CI 8--24\%)}
\newcommand{\StaticDiag}{0}         % of 11 vuln, overflow diagnosed by static stack
\newcommand{\StaticZero}{8}         % of 11, zero static diagnostics at all
\newcommand{\StaticOther}{3}        % of 11, unrelated (non-overflow) warning
\newcommand{\Nintraonly}{0}         % ASan-invisible intra-object cases in generated sample

\lstset{basicstyle=\ttfamily\small, breaklines=true, columns=fullflexible,
        keywordstyle=\bfseries, commentstyle=\itshape\color{gray},
        showstringspaces=false, frame=single, framesep=4pt, language=C}

\titleformat{\section}{\normalfont\large\bfseries}{\thesection}{0.6em}{}
\titleformat{\subsection}{\normalfont\normalsize\bfseries}{\thesubsection}{0.6em}{}

\title{\bfseries Detection Gaps for Spatial Memory Corruption in\\
LLM-Generated C: An Intra-Object Overflow Study Across\\
Detectors, Toolchains, and Code Models}
\author{Sadik Erisen\\\small \texttt{elmanoukdesign@gmail.com}}
\date{}

\begin{document}
\maketitle

\begin{abstract}
Large language models are increasingly used to generate C code that is compiled and run
inside sandboxes we trust to contain untrusted behaviour. That trust rests on a stack of
commodity defences: compiler warnings, static analyzers, build-time hardening, and
sanitizers. We ask a narrow, falsifiable question: for the specific class of spatial
memory corruption, which of these defences actually fire, and does the code a model
produces land in the gaps they leave?

We run three experiments. \textbf{H1} points ten independent commodity detectors at six
memory-corruption classes and records which are caught. One class ---
\emph{intra-object} (adjacent-field) overflow, an unbounded write that overwrites a
neighbouring struct field but stays inside a single valid allocation --- is missed by
\emph{every} detector, including AddressSanitizer and Valgrind, because no allocation
boundary is crossed. \textbf{H2} shows this gap is not a toolchain artefact: it holds
across five compilers (gcc-11, gcc-12, clang-14, clang-15, tcc), two \texttt{\_FORTIFY\_SOURCE}
levels, and AddressSanitizer under both gcc and clang, and across two CPU architectures.
\textbf{H3} asks whether real code models produce this class. We give three local code
models \Ntotal{} security-neutral tasks, each with a fixed-size struct field, and check
every generated snippet against a single-translation-unit, AddressSanitizer-backed ground
truth calibrated on golden-safe and golden-vulnerable references. \Nvuln{} of \Nusable{}
usable snippets contain a confirmed spatial overflow (a lower bound), concentrated almost
entirely in the smallest model (\RateQwenOneFive{} versus 0--1.4\% for the 7B models). The
commodity static stack (gcc \texttt{-Wall}, \texttt{-fanalyzer}, cppcheck) diagnosed the
overflow in \StaticDiag{} of the \Nvuln{}: it is structurally blind to input-length-dependent
overflow in a function it analyzes in isolation. We connect the three results at the
detection-layer level rather than by asserting more than we measured: the shipped
\emph{static} layer catches none of these generated overflows, and H1/H2 show the shipped
\emph{runtime} tool (AddressSanitizer) is blind to the intra-object subclass, so neither
layer is a sufficient detector for this class. One generated overflow (\texttt{parse\_kv})
writes \emph{through} an adjacent struct field before reaching the allocation boundary; had
the model bounded the loop, the same write would be an intra-object overflow invisible to
ASan --- the H1 probe in generated form. We report every number with its limits, including
why the rate is a lower bound and why a hand-verification pass removed one oracle false
positive.
\end{abstract}

\section{Introduction}
When an LLM-driven agent compiles and runs its own C code inside a sandbox, the sandbox
has to hold in two places at once: at the container boundary, and at the memory boundary
inside the process. This paper is only about the second. We do not argue about whether an
agent is ``capable'' of writing an exploit; we measure which memory-safety weaknesses
survive the commodity tooling we normally trust, and whether generated code falls into
the surviving classes.

The contribution is three linked, falsifiable measurements rather than a new tool:
\begin{itemize}[leftmargin=1.4em]
  \item \textbf{H1 (detector coverage).} Across ten commodity detectors and six
  memory-corruption classes, intra-object overflow is a \emph{total} detection gap ---
  missed by every static analyzer, every compiler warning set, build-time hardening, and
  both major sanitizers. We define the operational states (silent, detected, crashed,
  prevented, clean) precisely so ``missed'' is not a matter of opinion.
  \item \textbf{H2 (toolchain generality).} The gap is structural, not a property of one
  compiler version: it survives five compilers, two hardening levels, both sanitizer
  back-ends, and two architectures (aarch64 and x86\_64).
  \item \textbf{H3 (does generated code hit the layer that is blind).} Three local code
  models, given fixed-buffer tasks with no security framing, produce confirmed spatial
  overflows at a measurable rate --- concentrated in the smallest model --- and the shipped
  static layer diagnosed \StaticDiag{} of them. Ground truth is automated, calibrated, and
  then hand-verified; the verification pass removed one oracle false positive and confirmed
  the numbers rest on AddressSanitizer, not on inspection.
\end{itemize}

We are deliberate about what each number is and is not. The H3 vulnerable rate is a lower
bound; it is conditional on tasks that involve a fixed-size buffer; and it is measured on
three small open-weight models on one machine. Section~\ref{sec:threats} states every such
limit rather than burying it.

\section{Threat model and definitions}
\label{sec:defs}
\subsection{What ``the sandbox'' means here}
By sandbox we mean the stack normally trusted to hold the line around
compiled-and-executed code: the compiler and its hardening flags, the static analyzers
run in CI, and the sanitizers used in testing. We treat each as an independent
\emph{detector} and ask, per memory-corruption class, whether it fires.

\subsection{Operational states of a run}
For a build that exercises a bug on a given input, we classify the outcome by direct
evidence (the process's exit signal, the diagnostic it emits, and a self-check the probe
prints):
\begin{description}[leftmargin=1.4em, style=nextline]
  \item[silent] Corruption confirmed by the probe self-check, the process exits 0, and no
  defence emitted any diagnostic. Exploitable and undetected.
  \item[detected] A defence named the memory error and stopped the process
  (\texttt{\_FORTIFY\_SOURCE} \texttt{\_\_chk}, stack canary, a sanitizer report).
  \item[crashed] A raw \texttt{SIGSEGV}/\texttt{SIGBUS}/\texttt{SIGABRT} with no
  diagnostic: the corruption faulted (a denial-of-service outcome), but nothing explained
  it.
  \item[prevented] The unsafe operation was refused before corruption (an allocator
  returned \texttt{NULL} or aborted at the call), and the self-check reports clean.
  \item[clean] The bug did not take effect on this input.
\end{description}
The distinction that matters for this paper is \textbf{silent} versus everything else. A
\textbf{silent} outcome is the definition of a detection gap: real corruption, zero
signal.

\subsection{Intra-object versus boundary overflow}
An \emph{allocation-boundary} overflow writes past the end of a heap or stack object;
red-zone tools (AddressSanitizer) and shadow-memory tools (Valgrind) see it because the
write crosses into instrumented memory. An \emph{intra-object} overflow writes past one
struct field into an adjacent field of the \emph{same} allocation:
\begin{lstlisting}
struct rec { char name[16]; int uid; }; /* one allocation */
strcpy(r->name, attacker_input);        /* 20 bytes: name[16] overrun spills into uid */
\end{lstlisting}
No allocation boundary is crossed, so boundary-based tools stay silent by construction.
This is the class the paper is built around.

\section{H1: detector coverage of memory-corruption classes}
\label{sec:h1}
\subsection{Method}
We point ten independent commodity detectors at six memory-corruption classes and record
CAUGHT versus MISSED, saving raw per-cell evidence so every verdict is auditable. The
detectors are, static/compile-time: gcc \texttt{-Wall -Wextra}, gcc \texttt{-fanalyzer},
gcc LTO type-mismatch, clang \texttt{-Warray-bounds}, clang \texttt{--analyze}, cppcheck;
and runtime: \texttt{\_FORTIFY\_SOURCE=2}, AddressSanitizer, UndefinedBehaviorSanitizer,
Valgrind. Each runtime column is built with isolation flags so it measures exactly one
tool (canary and FORTIFY disabled when measuring a sanitizer, etc.). A compile diagnostic
counts only if a memory-safety checker fires on a real line; meta-noise (``unknown warning
option'') is filtered and never counts as a catch.

The six classes are ABI/type skew (CWE-843), out-of-bounds write in three variants
(intra-object, heap, stack; CWE-787), use-after-free (CWE-416), and integer-overflow
allocation (CWE-190).

\subsection{Results}
Table~\ref{tab:h1} gives the coverage. The headline: intra-object overflow is caught by
\emph{no} detector --- a total gap --- while the heap and stack variants of the same
CWE-787 are caught at runtime by ASan and Valgrind. Static analysis (all six compile-time
detectors) caught none of the six classes on these probes. Under the shipped build
(\texttt{-O2 -D\_FORTIFY\_SOURCE=2}, default canary), five of six classes are
\textbf{silent} and only the stack overflow is \textbf{detected} (by the canary).

\begin{table}[h]\centering\small
\caption{H1: which commodity detector catches which class. ``static'' = all six
compile-time detectors. Shipped state is the \texttt{-O2 -D\_FORTIFY\_SOURCE=2} build.}
\label{tab:h1}
\begin{tabular}{lccccl}
\toprule
class & static & FORTIFY & ASan & Valgrind & shipped state \\
\midrule
oob\_intra (adjacent field) & miss & miss & \textbf{miss} & \textbf{miss} & silent \\
oob\_heap                   & miss & miss & CAUGHT & CAUGHT & silent \\
oob\_stack                  & miss & miss & CAUGHT & CAUGHT & detected (canary) \\
use-after-free              & miss & --   & CAUGHT & CAUGHT & silent \\
int-overflow alloc          & miss & --   & partial & partial & silent \\
abi/type skew               & miss & --   & --     & --     & silent \\
\bottomrule
\end{tabular}
\end{table}

The result is not that these tools are weak in general --- ASan and Valgrind catch the
heap and stack overflows cleanly. It is that the \emph{intra-object} variant is invisible
to all of them at once, and that the shipped hardening bundle leaves five of six classes
silent. The intra-object row is the load-bearing cell of this paper and is re-verified in
Section~\ref{sec:h2} across five compilers and two architectures; the integer-overflow and
ABI-skew rows summarise the recorded detector matrix and mark those partial/absent cells
conservatively.

\section{H2: the gap is a structural property, not a toolchain artefact}
\label{sec:h2}
\subsection{Method}
We re-run the six classes through five compilers (gcc-11, gcc-12, clang-14, clang-15,
tcc) for compile-time warnings, through \texttt{\_FORTIFY\_SOURCE} levels 2 and 3 for the
shipped state, and through AddressSanitizer built by both gcc-12 and clang-15. H2 holds if
intra-object stays a total gap and no compiler or level newly catches it.

\subsection{Results}
No compiler warned on intra-object at any level; the shipped state is silent under
\texttt{\_FORTIFY\_SOURCE=2} and \texttt{=3}; and ASan misses it under both gcc and clang.
The gap holds across every toolchain tested. Two honest side-findings: gcc-12 is stricter
than gcc-11 (it emits a compile-time use-after-free warning where gcc-11 is silent), and
the five-of-six-silent shipped pattern is stable across all compilers. Because these
experiments ran on x86\_64 while the companion sandbox study ran on aarch64 with gcc-15,
the intra-object gap is confirmed across two architectures and a wide gcc version range
(11, 12, 15). The result is consistent with its cause: a boundary-based tool cannot see a
write that never crosses a boundary, regardless of compiler.

\section{H3: do real code models produce this class, and does the static stack miss it?}
\label{sec:h3}
\subsection{Method}
We test three local open-weight code models --- \texttt{qwen2.5-coder:7b},
\texttt{deepseek-coder:6.7b}, and \texttt{qwen2.5-coder:1.5b} --- on 18 security-neutral
tasks. Each task asks for a struct with a fixed-size \texttt{char} field and a function
that fills it from variable-length input; no task mentions safety, overflow, or bounds. We
draw four samples per task per model (one at temperature 0, three at 0.6), for \Ntotal{}
generations.

\paragraph{Provided contract.} To let an automated oracle bind to the output, each prompt
specifies the exact struct and function signature (\texttt{void fn(struct S *s, const char
*input)}). This fixes the interface; it does \emph{not} hint at safety. What we measure is
whether, given that interface, the model bounds the copy.

\paragraph{Two-oracle ground truth.} Every snippet is judged by an automated oracle, not by
inspection. The driver is concatenated into the \emph{same translation unit} as the
snippet, so it binds to the model's own struct and function; any signature or layout
deviation becomes a compile error we count as \texttt{compile\_fail} rather than silent
undefined behaviour. Under AddressSanitizer we sweep input lengths $\{N{-}1, N, N{+}1,
N{+}8, N{+}128\}$ for a field of size $N$ and apply two oracles:
\begin{description}[leftmargin=1.4em, style=nextline]
  \item[Oracle A (ASan).] A grossly oversized input that overruns the whole allocation
  triggers a heap-buffer-overflow report.
  \item[Oracle B (intra-object).] The trailing scalar field is set to a sentinel; an
  off-by-one or small overflow that stays \emph{inside} the allocation clobbers the
  sentinel --- the exact class Oracle A (and ASan) cannot see. Oracle B is applied only to
  tasks with a trailing scalar and a field size that is a multiple of the scalar's
  alignment (no padding between field and sentinel).
\end{description}
A snippet is vulnerable if either oracle fires.

\paragraph{Calibration.} Before any generated code is judged, every one of the 18 drivers
is validated against a golden-safe reference (\texttt{strncpy} + explicit null) and a
golden-vulnerable reference (\texttt{strcpy}); the driver is trusted only if it reports the
safe reference clean and the vulnerable reference flagged. All 18 drivers passed.

\paragraph{Static-stack classification.} Independently, each snippet is run through the
commodity static stack (gcc \texttt{-Wall -Wextra}, gcc \texttt{-fanalyzer}, cppcheck) on
the snippet alone. The money metric is: of the oracle-confirmed vulnerable snippets, how
many did the static stack flag?

\subsection{Results}
Of \Ntotal{} generations, \Nusable{} compiled against the oracle; \Ncfail{} were
\texttt{compile\_fail} (contract deviations: one corrupted generation, one write to a
\texttt{const} pointer, one \texttt{typedef}'d struct that left \texttt{struct header}
undefined). We exclude all three from the denominator; on inspection the \texttt{typedef}
case contains an unbounded copy and would likely be vulnerable, so excluding it keeps the
rate a lower bound.

\Nvuln{} of \Nusable{} usable snippets contain an AddressSanitizer-confirmed spatial
overflow, an overall rate of \RateAll{} (95\% CI \RateAllCI{}). The result is dominated by
model size (Table~\ref{tab:h3}): the smallest model, \texttt{qwen2.5-coder:1.5b}, produced
the class in \RateQwenOneFive{}, while \texttt{deepseek-coder:6.7b} produced \RateDeepseek{}
and \texttt{qwen2.5-coder:7b} produced \RateQwenSeven{}. We state this as an observation on
three points along the model-size axis, not a scaling law; the intervals are wide.

\begin{table}[h]\centering\small
\caption{H3: confirmed spatial overflow rate by model (usable snippets), and how the
commodity static stack diagnosed the confirmed cases.}
\label{tab:h3}
\begin{tabular}{lccc}
\toprule
model & confirmed / usable & rate & 95\% CI \\
\midrule
qwen2.5-coder:1.5b  & 10/70 & 14.3\% & 8--24\% \\
deepseek-coder:6.7b & 1/71  & 1.4\%  & 0--8\% \\
qwen2.5-coder:7b    & 0/72  & 0.0\%  & 0--5\% \\
\midrule
\multicolumn{4}{l}{\emph{static-stack diagnosis of the \Nvuln{} confirmed overflows:}}\\
\multicolumn{4}{l}{\quad overflow diagnosed by any static tool: \StaticDiag{}/\Nvuln{}}\\
\multicolumn{4}{l}{\quad zero static diagnostics at all: \StaticZero{}/\Nvuln{}}\\
\multicolumn{4}{l}{\quad unrelated warning on the function (truncation, sign-compare): \StaticOther{}/\Nvuln{}}\\
\bottomrule
\end{tabular}
\end{table}

\paragraph{The static layer is blind to this class.} No member of the commodity static
stack diagnosed the actual overflow in \emph{any} of the \Nvuln{} confirmed cases
(\StaticDiag{}/\Nvuln{}). \StaticZero{} drew no diagnostic at all; the remaining \StaticOther{}
drew a warning on the function that points at a \emph{different} defect (a
\texttt{-Wstringop-truncation} about missing null-termination, two \texttt{-Wsign-compare})
but not at the overflow AddressSanitizer confirmed. This is not a claim that the tools are
weak in general --- gcc's object-size checks fire when a size is statically known. It is
that these overflows are input-length-dependent writes in a function the analyzer sees in
isolation, where the caller's input length is opaque; that is precisely the real-world
condition of a parser compiled in its own translation unit, and it is undecidable at
compile time. Only runtime AddressSanitizer caught the \Nvuln{}, and ASan is not part of a
shipped build.

\paragraph{Where this meets H1/H2.} We do not claim generated code produced an
ASan-invisible intra-object overflow: in this sample it did not (\Nintraonly{} such cases;
the one candidate our intra-object oracle flagged was a false positive --- a function that
legitimately assigns its trailing field --- removed by hand-verification). The connection is
at the layer level. The shipped \emph{static} layer caught none of the \Nvuln{} generated
overflows (this section); H1/H2 show the shipped-adjacent \emph{runtime} tool is blind to
the intra-object subclass. Neither layer is a sufficient detector for the class. One case
makes the boundary concrete: \texttt{qwen2.5-coder:1.5b}'s \texttt{parse\_kv} copies into
\texttt{key[16]} with an unbounded loop; its ASan report shows the write reaching ``0 bytes
to the right of [the] 64-byte region,'' i.e.\ it passes \emph{through} the adjacent
\texttt{val[48]} field (offsets 16--63 of the 64-byte struct) and faults only at the
allocation boundary. Had the model bounded that loop to the struct --- as other samples did
--- the same \texttt{key}$\rightarrow$\texttt{val} overwrite would be an intra-object write
that never crosses a boundary and that AddressSanitizer would not report: the H1 probe,
arrived at from generated code.

\paragraph{The two-oracle design and what it contributed.} Our ground truth combined an
allocation-boundary oracle (A, AddressSanitizer) with an intra-object sentinel oracle (B).
In this run \emph{all \Nvuln{} confirmed findings came from oracle A}; oracle B produced no
true positives and one false positive (the legitimate trailing-field assignment above),
which hand-verification caught. We report this plainly: the rate rests on ASan. Oracle B's
value here was as a check for the ASan-invisible subclass that, in this sample, did not
occur.

\section{Threats to validity}
\label{sec:threats}
\begin{itemize}[leftmargin=1.4em]
  \item \textbf{The H3 rate is a lower bound.} Oracle B cannot see a mid-struct
  \texttt{field[i]}$\rightarrow$\texttt{field[i+1]} overflow, the multi-\texttt{char}-field
  tasks are judged by Oracle A alone, and one of the three excluded
  \texttt{compile\_fail}s contains an unbounded copy that would likely be vulnerable if it
  compiled. Every one of these effects can only raise the true rate above what we report.
  \item \textbf{Oracle B produced a false positive, and the rate rests on ASan.} The
  intra-object sentinel oracle assumes the function never writes its trailing scalar; a
  generated function that legitimately assigns that field (\texttt{s->count = 1}) trips the
  sentinel without any overflow. Our golden-safe calibration reference did not write the
  trailing field, so calibration passed without exercising this case. We caught the single
  occurrence by hand-verification and excluded it; all \Nvuln{} confirmed findings rest on
  AddressSanitizer (oracle A), not on oracle B. A corrected oracle B would compare the
  clobbered bytes against a partial-overflow signature (high bytes of the sentinel intact)
  rather than exact inequality.
  \item \textbf{Conditional on fixed-buffer tasks.} The tasks were chosen to involve a
  fixed-size buffer, so the rate is ``conditional on a fixed-buffer task,'' not a rate over
  all code a model writes.
  \item \textbf{Three small open models, one machine.} We test three open-weight models in
  the 1.5--7B range locally; we do not claim these rates transfer to frontier hosted
  models. The provided-contract design fixes the interface and may differ from free-form
  generation.
  \item \textbf{Standard tool configuration.} A non-default ASan build with intra-object
  red zones (\texttt{-fsanitize-address-\allowbreak field-padding}) can catch some
  intra-object overflows, but it is C++-oriented, opt-in, and part of no shipped build; the
  practical gap stands. We state this rather than omit it.
  \item \textbf{Compile-time versus runtime.} The static-stack-miss metric is a
  \emph{compile-time} claim and is kept distinct from the shipped-\texttt{\_FORTIFY\_SOURCE}
  runtime state; we do not blur a static miss into a runtime catch.
  \item \textbf{Single toolchain per version.} One gcc/clang/valgrind build each. The
  intra-object result is structural and H2 confirms it across two architectures, but the
  broader matrix should be re-run elsewhere to confirm.
\end{itemize}

\section{Related work}
AddressSanitizer~\cite{asan} and Valgrind/Memcheck~\cite{valgrind} are the red-zone and
shadow-memory tools whose boundary model is precisely why intra-object overflow is
invisible to them. The closest prior art on the intra-object case is ASan's own
field-padding extension~\cite{asanintra}, which inserts poisoned padding between fields to
catch adjacent-field overflow --- but only for C++ non-standard-layout classes, behind the
experimental \texttt{-fsanitize-address-\allowbreak field-padding} flag, and in no shipped
build; that scoping is exactly why the practical gap we measure stands.
\texttt{\_FORTIFY\_SOURCE}~\cite{fortify} provides the build-time \texttt{\_\_builtin\_\-object\_size}
checks we measure at the shipped-hardening layer. A separate line of work reports that code
models generate insecure code --- Pearce et al.~\cite{copilot} found roughly 40\% of
Copilot completions across 89 security-relevant scenarios were vulnerable. Our contribution
is not that claim but its connection to a measured, toolchain-general detector gap: we show
generated code lands in a spatial-overflow class the shipped detection stack cannot see, and
we separate the static layer (blind to input-dependent overflow in isolated functions) from
the runtime layer (blind to the intra-object subclass).

\section{Conclusion}
Spatial memory corruption has a specific, reproducible blind spot --- intra-object
overflow --- that no commodity detector catches, and that survives five compilers, two
hardening levels, both sanitizer back-ends, and two architectures. Independently, small
code models produce spatial overflows on ordinary fixed-buffer tasks at a measurable rate
(dominated by the smallest model), and the shipped static layer diagnosed none of them.
Read together at the layer level: the static layer misses input-length-dependent overflow
in isolated functions, and the runtime layer (AddressSanitizer) misses the intra-object
subclass, so a pipeline that compiles and runs LLM-generated C cannot rely on compiler
warnings, static analysis, build hardening, \emph{or} AddressSanitizer to be a sufficient
detector for this class. Detecting adjacent-field overflow needs an oracle that reasons
about struct-field boundaries, not allocation boundaries --- and that oracle is not in the
default stack today.

\paragraph{Artifacts.} The detector matrix, toolchain sweep, H3 generation-and-oracle
harness, all generated snippets, the calibration references, and raw per-cell evidence are
released with the paper.

\begin{thebibliography}{9}
\bibitem{asan} K.~Serebryany, D.~Bruening, A.~Potapenko, and D.~Vyukov.
  ``AddressSanitizer: A Fast Address Sanity Checker.''
  \emph{2012 USENIX Annual Technical Conference (USENIX ATC '12)}, pp.~309--318, 2012.
\bibitem{valgrind} N.~Nethercote and J.~Seward.
  ``Valgrind: A Framework for Heavyweight Dynamic Binary Instrumentation.''
  \emph{Proc.\ 28th ACM SIGPLAN Conf.\ on Programming Language Design and Implementation
  (PLDI '07)}, pp.~89--100, 2007.
\bibitem{fortify} J.~Jelinek.
  ``Object size checking to prevent (some) buffer overflows.''
  GCC patches mailing list, Sept.\ 2004.
  \url{https://gcc.gnu.org/legacy-ml/gcc-patches/2004-09/msg02055.html}.
  (Origin of \texttt{\_FORTIFY\_SOURCE}; see also the glibc manual.)
\bibitem{asanintra} Google Sanitizers project.
  ``AddressSanitizerIntraObjectOverflow.''
  \url{https://github.com/google/sanitizers/wiki/AddressSanitizerIntraObjectOverflow}.
  (Experimental \texttt{-fsanitize-address-field-padding}; C++ non-standard-layout classes.)
\bibitem{copilot} H.~Pearce, B.~Ahmad, B.~Tan, B.~Dolan-Gavitt, and R.~Karri.
  ``Asleep at the Keyboard? Assessing the Security of GitHub Copilot's Code Contributions.''
  \emph{2022 IEEE Symposium on Security and Privacy (SP)}, pp.~754--768, 2022.
\end{thebibliography}

\end{document}
